% =============================================================================
%  Relatório Técnico — Criptanálise Diferencial em Cifras de Fluxo
%  Baseado em: Biham & Dunkelman, "Differential Cryptanalysis in Stream Ciphers"
%  Implementação: LCG sobre GF(2^e) com função filtro de baixo grau algébrico
% =============================================================================
\documentclass[12pt,a4paper]{article}

% ─── Pacotes ─────────────────────────────────────────────────────────────────
\usepackage[utf8]{inputenc}
\usepackage[T1]{fontenc}
\usepackage[brazilian]{babel}
\usepackage{amsmath, amssymb, amsthm}
\usepackage{mathtools}
\usepackage{geometry}
\usepackage{hyperref}
\usepackage{xcolor}
\usepackage{listings}
\usepackage{booktabs}
\usepackage{array}
\usepackage{multirow}
\usepackage{enumitem}
\usepackage{tcolorbox}
\usepackage{fancyhdr}
\usepackage{titlesec}
\usepackage{graphicx}
\usepackage{float}
\usepackage{bm}

\geometry{
  left=2.5cm, right=2.5cm,
  top=3cm, bottom=3cm
}

% ─── Cores e estilos ─────────────────────────────────────────────────────────
\definecolor{deepblue}{RGB}{0,0,153}
\definecolor{codegreen}{RGB}{0,128,0}
\definecolor{codegray}{RGB}{128,128,128}
\definecolor{backcolor}{RGB}{245,245,245}
\definecolor{theoremblue}{RGB}{230,240,255}
\definecolor{defgreen}{RGB}{230,255,230}
\definecolor{warnred}{RGB}{255,230,230}
\definecolor{darkred}{RGB}{150,0,0}

\hypersetup{
  colorlinks=true,
  linkcolor=deepblue,
  urlcolor=deepblue,
  citecolor=darkred
}

% ─── Ambientes matemáticos ───────────────────────────────────────────────────
\newtheorem{theorem}{Teorema}[section]
\newtheorem{lemma}[theorem]{Lema}
\newtheorem{corollary}[theorem]{Corolário}
\newtheorem{proposition}[theorem]{Proposição}
\theoremstyle{definition}
\newtheorem{definition}[theorem]{Definição}
\newtheorem{example}[theorem]{Exemplo}
\theoremstyle{remark}
\newtheorem{remark}[theorem]{Observação}

% ─── Caixa de referência ao paper ────────────────────────────────────────────
\newtcolorbox{paperref}{
  colback=theoremblue,
  colframe=deepblue,
  fonttitle=\bfseries,
  title={Referência ao Paper (Biham \& Dunkelman)},
  left=4pt, right=4pt, top=4pt, bottom=4pt
}

\newtcolorbox{mathbox}{
  colback=defgreen,
  colframe=codegreen!70!black,
  fonttitle=\bfseries,
  left=4pt, right=4pt, top=4pt, bottom=4pt
}

\newtcolorbox{warnbox}{
  colback=warnred,
  colframe=darkred,
  fonttitle=\bfseries,
  title={Limitação de Segurança},
  left=4pt, right=4pt, top=4pt, bottom=4pt
}

% ─── Cabeçalho/rodapé ────────────────────────────────────────────────────────
\pagestyle{fancy}
\fancyhf{}
\fancyhead[L]{\small\textit{Criptanálise Diferencial em Cifras de Fluxo}}
\fancyhead[R]{\small\textit{LCG sobre GF$(2^e)$ com Função Filtro}}
\fancyfoot[C]{\thepage}

% ─── Notação ─────────────────────────────────────────────────────────────────
\newcommand{\GF}[1]{\mathbb{F}_{2^{#1}}}
\newcommand{\GFtwo}{\mathbb{F}_2}
\newcommand{\xor}{\oplus}
\newcommand{\AND}{\wedge}
\newcommand{\Diff}[2]{D^{#1}_{#2}}
\newcommand{\filter}{f}
\newcommand{\state}{x}
\newcommand{\ks}{z}
\newcommand{\lcgmap}[2]{#1 \cdot #2 + c}
\newcommand{\bitk}[2]{s_{#1}(#2)}
\newcommand{\mat}[1]{\mathbf{#1}}

% =============================================================================
\begin{document}

% ─── Capa ────────────────────────────────────────────────────────────────────
\begin{titlepage}
  \centering
  \vspace*{2cm}
  {\Huge\bfseries Criptanálise Diferencial\\[0.4em]
   em Cifras de Fluxo\par}
  \vspace{1cm}
  {\Large\itshape Análise Matemática e Implementação de Ataques\\
   sobre LCG em $\GF{e}$ com Função Filtro de Baixo Grau\par}
  \vspace{2cm}
  \rule{0.8\textwidth}{0.4pt}\\[0.5em]
  {\large
    Baseado em:\\[0.3em]
    \textbf{Eli Biham \& Orr Dunkelman}\\
    \textit{``Differential Cryptanalysis in Stream Ciphers''}\\
    \textit{(Biham \& Dunkelman, 2007)}
  }\\[0.5em]
  \rule{0.8\textwidth}{0.4pt}
  \vspace{2cm}

  {\large Matheus Sousa\\
    Departamento de Eletrônica e Sistemas --- UFPE}\\[0.5em]
  {\large 9 de março de 2026}
  \vfill
  {\small
    \textbf{Resumo.} Este relatório apresenta a fundamentação matemática
    completa e a análise da implementação dos ataques de criptanálise diferencial
    descritos por Biham e Dunkelman aplicados a uma cifra de fluxo sintética
    modelada como um \emph{Gerador Congruencial Linear} (LCG) sobre o corpo
    finito $\GF{e}$, combinado com uma função filtro booleana de grau algébrico
    baixo. São provados, com álgebra booleana e linear, todos os invariantes
    explorados pelos seis módulos de ataque do projeto:
    recuperação de chave por sistema linear, correlação diferencial, ataque
    distinguidor, diferencial de segunda ordem, múltiplas características e
    aceleração da busca exaustiva. Por fim, analisa-se a extensão para funções
    filtro de grau algébrico arbitrário $d$, explicitando o custo $2^d$ do
    diferencial de ordem $d$ e as condições de inviabilidade computacional.
  }
\end{titlepage}

\tableofcontents
\newpage

% =============================================================================
\section{Introdução e Modelo do Sistema}
% =============================================================================

\subsection{Contextualização}

\begin{paperref}
  Biham \& Dunkelman definem no \S3.1 as três categorias de características
  diferenciais para cifras síncronas:
  \[
    (\Delta_{\text{key}},\Delta_{IV}) \xrightarrow{\text{INIT}} \Delta S,
    \qquad
    \Delta S \xrightarrow{\text{UPDATE}} \Delta S,
    \qquad
    \Delta S \xrightarrow{\text{OUTPUT}} \Delta KS.
  \]
  O projeto implementa exatamente este modelo, explorando a composição
  $\Delta_0 \xrightarrow{\text{INIT}+\text{OUTPUT}} \Delta KS$ com
  probabilidade 1 (característica perfeita).
\end{paperref}

A cifra de fluxo estudada neste projeto é formada por dois componentes:

\begin{enumerate}[label=(\roman*)]
  \item \textbf{Gerador de estado interno (UPDATE):} um Gerador Congruencial
        Linear (LCG) definido sobre o corpo finito $\GF{e}$:
        \[
          x_{t+1} = a \cdot x_t + c, \qquad a,c,x_0 \in \GF{e},
        \]
        onde $\GF{e} \cong \GFtwo[x] / \langle h(x) \rangle$, com $h(x)$
        polinômio primitivo de grau $e$ sobre $\GFtwo$.

  \item \textbf{Função filtro (OUTPUT):} uma função booleana
        $\filter : \GF{e} \to \GFtwo$ que extrai 1 bit de keystream por clock:
        \[
          \ks_t = \filter(x_t) = s_{20}(x_t) \xor \bigl(s_0(x_t) \AND s_{10}(x_t)\bigr),
        \]
        onde $s_k(x)$ denota o coeficiente do monômio $x^k$ na representação
        polinomial de $x$ em $\GF{e}$.
\end{enumerate}

\subsection{Representação em $\GF{e}$}

Cada estado $x_t \in \GF{e}$ é representado como um vetor de $e$ bits:
\[
  x_t = \sum_{k=0}^{e-1} s_k(x_t)\, \xi^k, \qquad s_k(x_t) \in \GFtwo,
\]
onde $\xi$ é raiz de $h$. A multiplicação por $a$ em $\GF{e}$ é uma
transformação $\GFtwo$-linear, fato fundamental para a validade de todos
os ataques.

\subsection{Configurações de Teste}

O arquivo \texttt{unified\_config.json} fornece 8 instâncias cobrindo
diferentes tamanhos de campo:

\begin{table}[H]
  \centering
  \caption{Configurações de campo do projeto}
  \begin{tabular}{clll}
    \toprule
    \textbf{Config} & \textbf{Campo} & \textbf{$h(x)$} & \textbf{$a(x)$} \\
    \midrule
    1 & $\GF{31}$  & $x^{31}+x^{13}+x^8+x^3+1$ & $x^4+1$ \\
    2 & $\GF{64}$  & $x^{64}+x^{11}+x^2+x+1$   & $x^4+1$ \\
    3 & $\GF{64}$  & $x^{64}+x^{42}+x^{33}+x^{22}+1$ & $x^5+1$ \\
    4 & $\GF{128}$ & $x^{128}+x^{79}+x^{51}+x^{29}+1$ & $x^5+1$ \\
    5 & $\GF{256}$ & $x^{256}+x^{213}+x^{141}+x^{51}+1$ & $x^6+1$ \\
    6 & $\GF{271}$ & $x^{271}+x^{188}+x^{121}+x^{93}+1$ & (pentanômio) \\
    7 & $\GF{291}$ & $x^{291}+x^{240}+x^{202}+x^{178}+1$ & (pentanômio) \\
    8 & $\GF{451}$ & $x^{451}+x^{402}+x^{292}+x^9+1$    & (pentanômio) \\
    \bottomrule
  \end{tabular}
\end{table}

% =============================================================================
\section{Fundamentação: Linearidade do LCG e Propagação de Diferenças}
% =============================================================================

\begin{paperref}
  Biham \& Dunkelman (\S4, parágrafo inicial): \textit{``It is easy to see
  that given two (key,IV) pairs \ldots the difference in the internal state
  is known for any given number of clocks. This allows to estimate the output
  difference of the combining function.''}  A prova formal abaixo estabelece
  este resultado para o nosso LCG.
\end{paperref}

\begin{theorem}[Propagação Linear de Diferenças no LCG]
  \label{thm:linear-prop}
  Sejam $x_0, x_0' \in \GF{e}$ dois estados iniciais com diferença
  $\Delta_0 = x_0 \xor x_0'$. Após $t$ clocks do LCG $x \mapsto ax+c$,
  a diferença entre os estados é:
  \[
    \Delta_t \coloneqq x_t \xor x_t' = a^t \cdot \Delta_0.
  \]
\end{theorem}

\begin{proof}
  Por indução em $t$. \textbf{Base:} $t=0$, $\Delta_0 = x_0 \xor x_0'$.\;
  \textbf{Passo:} supondo $x_t \xor x_t' = a^t \Delta_0$,
  \begin{align*}
    x_{t+1} \xor x_{t+1}' &= (a x_t + c) \xor (a x_t' + c) \\
                           &= a(x_t \xor x_t') \xor \underbrace{(c \xor c)}_{0}
                              \quad \text{(linearidade de } \GF{e}\text{)} \\
                           &= a \cdot a^t \Delta_0 = a^{t+1} \Delta_0. \qedhere
  \end{align*}
\end{proof}

\begin{remark}
  O termo $c$ cancela exatamente porque a subtração em $\GFtwo$ coincide com
  a adição ($x \xor x = 0$). Isso é a causa central da vulnerabilidade: a
  diferença evolui \emph{deterministicamente} e é completamente previsível pelo
  atacante a partir do momento em que $\Delta_0$ é fixado.
\end{remark}

% =============================================================================
\section{Módulo 1: \texttt{batch\_attack} + \texttt{recover\_key}}
\label{sec:batch}
% =============================================================================

\subsection{Estratégia do Ataque}

\begin{paperref}
  Biham \& Dunkelman (\S4.1.3, \textit{Differential Key Recovery Attack on
  Toyocrypt}): \textit{``Given two key streams produced by two keys whose
  initial state difference is $\Delta S = s_0$, it is possible to reconstruct
  most of the internal state.''} No projeto, a injeção de $\Delta_0 = 1$
  (i.e., $s^0 = x^0$, o monômio constante em $\GF{e}$) serve ao mesmo
  propósito.
\end{paperref}

O atacante fixa $\Delta_0 = 1 \in \GF{e}$ (bit 0 invertido) e observa,
para cada $t$, o par $(z_t, z_t')$ onde $z_t = \filter(x_t)$ e
$z_t' = \filter(x_t \xor \Delta_0)$ (sonda local, não requer dois fluxos
simultâneos).

\subsection{Prova: Linearização da Diferença de Saída}

\begin{theorem}[Extração de bit de estado via diferença de saída]
  \label{thm:delta-z-s10}
  Seja $\filter(x) = s_{20}(x) \xor (s_0(x) \AND s_{10}(x))$ e
  $\Delta_0 = 1$ (monômio $x^0$). Para qualquer estado $x_t \in \GF{e}$:
  \[
    \Delta \ks_t \coloneqq \filter(x_t) \xor \filter(x_t \xor \Delta_0) = s_{10}(x_t).
  \]
\end{theorem}

\begin{proof}
  Seja $p = x_t$ e $p' = x_t \xor 1$. Então $s_k(p') = s_k(p)$ para todo
  $k \neq 0$, e $s_0(p') = 1 \xor s_0(p)$. Portanto:
  \begin{align*}
    \filter(p) &= s_{20}(p) \xor (s_0(p) \AND s_{10}(p)), \\
    \filter(p') &= s_{20}(p) \xor \bigl((s_0(p) \xor 1) \AND s_{10}(p)\bigr).
  \end{align*}
  Calculando $\filter(p) \xor \filter(p')$:
  \begin{align*}
    \Delta\ks &= \Bigl[s_{20} \xor (s_0 \AND s_{10})\Bigr]
               \xor \Bigl[s_{20} \xor ((s_0 \xor 1) \AND s_{10})\Bigr] \\
              &= (s_0 \AND s_{10}) \xor \bigl[(s_0 \xor 1) \AND s_{10}\bigr] \\
              &= s_{10} \AND \bigl[s_0 \xor (s_0 \xor 1)\bigr]
               \quad\text{(distributividade de AND sobre XOR)} \\
              &= s_{10} \AND 1 = s_{10}. \qedhere
  \end{align*}
\end{proof}

\begin{remark}
  A prova usa a identidade distributiva da álgebra de Boole:
  $(A \xor B) \AND C = (A \AND C) \xor (B \AND C)$, e o fato de que
  $s_0 \xor (s_0 \xor 1) = 1$ em $\GFtwo$.
\end{remark}

\subsection{Prova: Validade no Tempo (Propagação pelo LCG)}

\begin{theorem}[Extração temporal de $s_{10}$]
  \label{thm:temporal}
  A sonda local $\filter(x_t) \xor \filter(x_t \xor \Delta_0)$ extrai
  $s_{10}(x_t)$ com probabilidade 1 para \emph{todo} $t \geq 0$.
\end{theorem}

\begin{proof}
  Em cada clock o atacante computa $x_t \xor \Delta_0$ \emph{diretamente},
  sem precisar propagar $\Delta_0$; basta inverter o bit 0 do estado corrente.
  Pelo Teorema~\ref{thm:delta-z-s10}, o resultado $\Delta\ks_t = s_{10}(x_t)$
  vale para qualquer $x_t$, em particular para $x_t = a^t x_0 + c_{t}$
  onde $c_t = \sum_{i=0}^{t-1} a^i c$ é a soma de deslocamento acumulada.
\end{proof}

\subsection{Recuperação do Estado Inicial: Sistema Linear}

\begin{paperref}
  O método é o análogo algébrico da recuperação descrita por Biham \&
  Dunkelman (\S4.1.3) para a Toyocrypt: identificados $e$ bits de estado
  ao longo do tempo, monta-se um sistema linear invertível.
\end{paperref}

Após observar $e = \deg(h)$ clocks, o atacante dispõe do vetor
$B = (\Delta\ks_0, \ldots, \Delta\ks_{e-1})^T = (s_{10}(x_0), \ldots,
s_{10}(x_{e-1}))^T \in \GFtwo^e$.

Escrevemos $x_t = a^t x_0 + C_t$ onde $C_t = \sum_{i=0}^{t-1}a^i c$.
O bit $s_{10}(x_t)$ é a projeção linear $\langle e_{10}, x_t \rangle$ no
sentido de $\GFtwo^e$, portanto:
\[
  B_t = s_{10}(a^t x_0 + C_t) = s_{10}(a^t x_0) \xor s_{10}(C_t).
\]

Definindo $Y_t \coloneqq B_t \xor s_{10}(C_t)$, temos $Y_t = s_{10}(a^t x_0)$.
O mapeamento $x_0 \mapsto s_{10}(a^t x_0)$ é $\GFtwo$-linear; sua matriz
$M \in \GFtwo^{e \times e}$ tem entradas:
\[
  M_{t,j} = s_{10}(a^t \cdot \xi^j), \qquad t,j \in \{0,\ldots,e-1\},
\]
onde $\xi^j$ é o $j$-ésimo elemento da base canônica de $\GF{e}$.

\begin{theorem}[Invertibilidade da Matriz de Recuperação]
  \label{thm:invertible}
  Se $a$ é um elemento primitivo de $\GF{e}$ (ou, mais geralmente, se
  o sequenciador $x \mapsto ax+c$ tem período $2^e - 1$), então $M$ é
  invertível sobre $\GFtwo$ e o estado $x_0$ é unicamente recuperado por:
  \[
    \vec{S}_0 = M^{-1} Y,
  \]
  onde $\vec{S}_0 \in \GFtwo^e$ é o vetor de bits do estado inicial $x_0$.
\end{theorem}

\begin{proof}[Prova (esboço)]
  As $e$ funções lineares $\ell_t(x_0) = s_{10}(a^t x_0)$ são
  $\GFtwo$-linearmente independentes porque o mapa $t \mapsto a^t$ percorre
  uma órbita de comprimento $2^e - 1$ no grupo multiplicativo $\GF{e}^*$,
  e a projeção $s_{10}$ é uma forma linear não-nula em $\GF{e}$.
  A linearidade independente das linhas de $M$ garante $\det(M) \neq 0$
  sobre $\GFtwo$.
\end{proof}

\subsection{Correspondência com o Código}

Em \texttt{recover\_key.cpp}, o laço constrói exatamente $M$ e $Y$:
\begin{itemize}
  \item Linha \texttt{M[i][j] = IsOne(coeff(rep(a\_pow * xj), target\_bit))}
        $\leftrightarrow M_{i,j} = s_{10}(a^i \cdot \xi^j)$.
  \item Linha \texttt{Y[i] = B\_i \^{} c\_bit} $\leftrightarrow Y_i = B_i \xor s_{10}(C_i)$.
  \item \texttt{inv(M\_inv, M)} resolve $M^{-1}$ em $\GFtwo^{e\times e}$ pela biblioteca NTL.
\end{itemize}

% =============================================================================
\section{Módulo 2: \texttt{differential\_correlation}}
\label{sec:corr}
% =============================================================================

\begin{paperref}
  Biham \& Dunkelman (\S4, segundo parágrafo): \textit{``This fact can also
  be used to distinguish the output of the stream cipher under two related
  keys (or IVs) from two random strings. Moreover, in some cases, it can be
  used to identify information about the internal state.''}
  O coeficiente de correlação $\phi$ de Matthews quantifica precisamente o
  viés que Biham \& Dunkelman descrevem informalmente.
\end{paperref}

\subsection{Correlação Passiva: $\phi(z_t, s_{20})$}

A função $\filter(x) = s_{20} \xor (s_0 \AND s_{10})$ possui um viés
mensurável entre a saída $\ks_t$ e o bit $s_{20}(x_t)$:

\begin{proposition}[Viés passivo da saída]
  \label{prop:passive-bias}
  Assumindo $x_t$ uniformemente distribuído em $\GF{e}$, a correlação
  de Matthews entre $\ks_t = \filter(x_t)$ e $s_{20}(x_t)$ é:
  \[
    \phi(\ks_t, s_{20}(x_t)) > 0.
  \]
\end{proposition}

\begin{proof}
  Calculamos a tabela de contingência. Como $s_0, s_{10}, s_{20}$ são bits
  estatisticamente independentes e uniformes em $\{0,1\}$, temos
  $\Pr[s_0=1]=\Pr[s_{10}=1]=\Pr[s_{20}=1]=1/2$.

  Para $\phi(\ks, s_{20}) \neq 0$, basta mostrar $P(\ks=s_{20}) \neq 1/2$:
  \begin{align*}
    P(\ks = s_{20}) &= P\bigl[s_{20} \xor (s_0 \AND s_{10}) = s_{20}\bigr] \\
                    &= P\bigl[s_0 \AND s_{10} = 0\bigr] \\
                    &= 1 - P[s_0=1] \cdot P[s_{10}=1] = 1 - \tfrac{1}{4} = \tfrac{3}{4}.
  \end{align*}
  Portanto $\phi > 0$ (correlação positiva). \qedhere
\end{proof}

\subsection{Correlação Diferencial: $\phi(\Delta\ks_t, s_{10})$}

\begin{theorem}[Correlação diferencial perfeita]
  \label{thm:diff-corr}
  Seja $\Delta\ks_t = \filter(x_t) \xor \filter(x_t \xor 1)$. Então:
  \[
    \phi(\Delta\ks_t,\, s_{10}(x_t)) = 1.
  \]
\end{theorem}

\begin{proof}
  Pelo Teorema~\ref{thm:delta-z-s10}, $\Delta\ks_t = s_{10}(x_t)$ com
  probabilidade 1, portanto a concordância é total: os eventos
  $\{\Delta\ks_t = 1\}$ e $\{s_{10}(x_t) = 1\}$ são idênticos.
  O coeficiente de Matthews de dois eventos idênticos é 1. \qedhere
\end{proof}

\subsection{Coeficiente de Matthews}

O código usa o coeficiente $\phi$ de Matthews definido como:
\[
  \phi = \frac{n_{11}n_{00} - n_{10}n_{01}}
         {\sqrt{(n_{11}+n_{10})(n_{01}+n_{00})(n_{11}+n_{01})(n_{10}+n_{00})}},
\]
onde $n_{ij}$ conta os pares em que a variável $X=i$ e $Y=j$.
Pelo Teorema~\ref{thm:diff-corr}, para o par $(\Delta\ks, s_{10})$:
$n_{10} = n_{01} = 0$, logo $\phi = 1$.

% =============================================================================
\section{Módulo 3: \texttt{distinguish\_attack}}
\label{sec:distinguish}
% =============================================================================

\begin{paperref}
  Biham \& Dunkelman (\S4.1.1, \textit{Distinguishing Attack on Toyocrypt}):
  \textit{``given two streams generated by related keys, we check whether the
  first output bit is the same and the second one differs. If this is the
  case, we output that the stream cipher in use is Toyocrypt.''}
  O distinguidor do projeto é uma generalização determinística desta ideia:
  identifica os clocks em que a diferença interna $\Delta_t$ tem $\Delta_0 =
  \Delta_{10} = 0$, momento em que $\Delta\ks_t = \Delta_{20}$ com
  probabilidade 1.
\end{paperref}

\subsection{Condição de Distinguibilidade}

\begin{theorem}[Distinguidor Determinístico]
  \label{thm:distinguisher}
  Seja $\Delta_t = a^t \Delta_0$ com $\Delta_0 = 1$. Para os clocks $t$ tais
  que $s_0(\Delta_t) = s_{10}(\Delta_t) = 0$, a diferença de saída satisfaz:
  \[
    \Delta\ks_t = s_{20}(\Delta_t).
  \]
  Portanto, o evento \emph{``par de streams com diferença inicial $\Delta_0=1$"} é
  distinguível de dois streams aleatórios independentes com taxa de erro 0\%.
\end{theorem}

\begin{proof}
  Seja $\Delta = \Delta_t$ no instante $t$, e sejam $p = x_t$, $p' = x_t \xor
  \Delta$. Por definição, $s_k(p') = s_k(p) \xor s_k(\Delta)$ para todo $k$
  (linearidade da projeção).

  Calculamos $\filter(p) \xor \filter(p')$:
  \begin{align*}
    \Delta\ks &= \bigl[s_{20}(p) \xor (s_0(p) \AND s_{10}(p))\bigr] \\
              &\quad \xor \bigl[(s_{20}(p)\xor\delta_{20}) \xor
                              ((s_0(p)\xor\delta_0) \AND (s_{10}(p)\xor\delta_{10}))
                              \bigr]
  \end{align*}
  onde $\delta_k = s_k(\Delta)$. Expandindo:
  \begin{align*}
    \Delta\ks &= \delta_{20} \xor
                 \bigl[s_0 \AND s_{10}\bigr] \xor
                 \bigl[(s_0\xor\delta_0) \AND (s_{10}\xor\delta_{10})\bigr].
  \end{align*}
  Expandindo o produto booleano:
  \[
    (s_0\xor\delta_0)(s_{10}\xor\delta_{10}) =
      s_0 s_{10} \xor s_0\delta_{10} \xor \delta_0 s_{10} \xor \delta_0\delta_{10}.
  \]
  Portanto:
  \[
    \Delta\ks = \delta_{20} \xor s_0\delta_{10} \xor \delta_0 s_{10}
                \xor \delta_0\delta_{10}.
  \]
  Sob a condição $\delta_0 = s_0(\Delta_t) = 0$ \textbf{e} $\delta_{10} =
  s_{10}(\Delta_t) = 0$:
  \[
    \Delta\ks = \delta_{20} \xor 0 \xor 0 \xor 0 = \delta_{20} = s_{20}(\Delta_t). \qedhere
  \]
\end{proof}

\subsection{Fórmula Geral da Diferença de Saída}

\begin{corollary}[Equação diferencial completa]
  \label{cor:full-diff}
  Para quaisquer $\delta_0, \delta_{10}, \delta_{20} \in \GFtwo$:
  \[
    \Delta\ks_t = \delta_{20} \xor (s_0(x_t) \AND \delta_{10})
                             \xor (\delta_0 \AND s_{10}(x_t))
                             \xor (\delta_0 \AND \delta_{10}).
  \]
\end{corollary}

Esta fórmula unificada é usada diretamente pelo \texttt{multi\_difference\_attack}
(Seção~\ref{sec:multi}) com diferentes valores de $\Delta_0$.

% =============================================================================
\section{Módulo 4: \texttt{higher\_order\_differential}}
\label{sec:higher2}
% =============================================================================

\begin{paperref}
  Biham \& Dunkelman (\S4, parágrafo sobre monômio de grau 3): \textit{``if
  there is a difference in only one bit which is combined through a monomial of
  degree 3, there is a probability of 1/4 that there is a difference in the
  key streams in the corresponding bit.''} O diferencial de ordem superior é
  exatamente a ferramenta para contornar esta probabilidade fracionária e
  recuperar informação determinística quando o grau algébrico é conhecido.
\end{paperref}

\subsection{Diferenciais de Ordem $k$: Definição}

\begin{definition}[Diferencial de ordem $k$]
  Dada $f : \GFtwo^n \to \GFtwo$ e vetores $\alpha_1, \ldots, \alpha_k \in
  \GFtwo^n$, o \emph{diferencial de ordem $k$} de $f$ no ponto $x$ é:
  \[
    D^k_{\alpha_1,\ldots,\alpha_k} f(x)
    \coloneqq \bigoplus_{S \subseteq [\![k]\!]} f\!\Bigl(x \xor
      \bigoplus_{i \in S} \alpha_i\Bigr),
  \]
  onde a soma é sobre todos os $2^k$ subconjuntos $S$ de $\{1,\ldots,k\}$.
\end{definition}

\subsection{Teorema Central: Grau Algébrico e Colapso do Diferencial}

\begin{theorem}[Colapso do diferencial de ordem $d$]
  \label{thm:collapse}
  Seja $f : \GFtwo^n \to \GFtwo$ uma função booleana de grau algébrico
  exato $d$ (i.e., $f$ possui pelo menos um monômio de grau $d$ em sua
  Forma Normal Algébrica (ANF), mas nenhum de grau $> d$). Então:
  \[
    D^d_{\alpha_1,\ldots,\alpha_d} f(x) = \text{const} \in \GFtwo
    \quad \forall\, x \in \GFtwo^n,
  \]
  e $D^{d-1} f$ não é constante em geral.
\end{theorem}

\begin{proof}[Prova por indução no grau]
  \textbf{Base ($d=1$):} $f(x) = \langle a, x \rangle \xor c$ (linear).
  $D^1_\alpha f(x) = f(x\xor\alpha) \xor f(x) = \langle a,\alpha\rangle$,
  constante. $\checkmark$

  \textbf{Passo:} Seja $f = m + g$ onde $m = x_{i_1}\cdots x_{i_d}$ é um
  monômio de grau $d$ e $g$ tem grau $< d$. Por hipótese de indução,
  $D^d g \equiv 0$ (constante 0, pois $g$ tem grau $< d$ e um diferencial
  de ordem maior que o grau anula a função). Resta calcular $D^d m$:
  \begin{align*}
    D^d_{\alpha_1,\ldots,\alpha_d} m(x)
    &= \bigoplus_{S \subseteq [\![d]\!]} \prod_{j=1}^d
       \Bigl(x_{i_j} \xor \bigoplus_{s \in S} \alpha_{s,i_j}\Bigr).
  \end{align*}
  Expandindo via inclusão-exclusão booleana sobre os $2^d$ termos,
  o único coeficiente de grau 0 (independente de $x$) é o produto
  das diferenças nos $d$ bits participantes:
  \[
    D^d_{\alpha_1,\ldots,\alpha_d} m(x)
    = \prod_{j=1}^d \Bigl(\bigoplus_{i: \alpha_i \text{ contribui ao bit } i_j}
      \alpha_{i,i_j}\Bigr) = \text{const}.
  \]
  Todos os termos que contêm variáveis $x_{i_j}$ cancelam em pares.
  Portanto $D^d f = D^d m \xor D^d g = \text{const} \xor 0 = \text{const}$.
\end{proof}

\subsection{Aplicação à Função Filtro $\filter$ (Grau 2)}

\begin{corollary}[Diferencial de 2ª ordem da função filtro]
  \label{cor:d2-filter}
  Para $\filter(x) = s_{20} \xor (s_0 \AND s_{10})$ e
  $\alpha = x^0$, $\beta = x^{10}$ (base canônica de $\GF{e}$):
  \[
    D^2_{\alpha,\beta} \filter(x) = \filter(x) \xor \filter(x\xor\alpha)
                       \xor \filter(x\xor\beta) \xor \filter(x\xor\alpha\xor\beta)
    = (a_0 \AND b_{10}) \xor (a_{10} \AND b_0),
  \]
  onde $a_k = s_k(\alpha)$, $b_k = s_k(\beta)$.
\end{corollary}

\begin{proof}
  Denotando $\delta_k = s_k(\Delta)$ para cada diferença $\Delta$, o
  Corolário~\ref{cor:full-diff} nos dá:
  \begin{align*}
    D^2_{\alpha,\beta} \filter(x)
    &= \Delta_\alpha \filter(x) \xor \Delta_\alpha \filter(x \xor \beta) \\
    &= \Bigl[a_{20} \xor s_0 a_{10} \xor a_0 s_{10} \xor a_0 a_{10}\Bigr] \\
    &\quad \xor \Bigl[a_{20} \xor (s_0 \xor b_0) a_{10}
                       \xor a_0 (s_{10}\xor b_{10}) \xor a_0 a_{10}\Bigr] \\
    &= (b_0 a_{10}) \xor (a_0 b_{10}),
  \end{align*}
  pois os termos com $s_0, s_{10}$ e $a_{20}$ e $a_0 a_{10}$ cancelam.
  Para $\alpha = x^0$ temos $a_0 = 1, a_{10} = 0$; para $\beta = x^{10}$
  temos $b_0 = 0, b_{10} = 1$; logo:
  \[
    D^2 = (0 \AND 0) \xor (1 \AND 1) = 0 \xor 1 = 1 = \text{const.} \qedhere
  \]
\end{proof}

A constância do diferencial de 2ª ordem, confirmada empiricamente em
\texttt{higher\_order\_differential.cpp} com taxa 1.0 em todas as 8
configurações, prova que $\filter$ tem grau exato 2.

% =============================================================================
\section{Módulo 5: \texttt{multi\_difference\_attack}}
\label{sec:multi}
% =============================================================================

\begin{paperref}
  Biham \& Dunkelman (\S4.1.3): \textit{``by disregarding the output
  difference when there is a difference in $s_{10}$, $s_{23}$, $s_{32}$,
  or $s_{42}$, almost the entire register can easily be recovered.''}
  No projeto, a mesma ideia é generalizada: múltiplas diferenças iniciais
  $\Delta^{(i)}$ geram características complementares que permitem recuperar
  bits específicos do estado.
\end{paperref}

\subsection{Recuperação de Bits por Características Complementares}

A partir do Corolário~\ref{cor:full-diff}:
\[
  \Delta\ks_t = \delta_{20} \xor (s_0 \AND \delta_{10})
               \xor (\delta_0 \AND s_{10}) \xor (\delta_0 \AND \delta_{10}).
\]

Dado $\Delta$ conhecido, o atacante pode isolar cada bit:

\begin{proposition}[Isolamento de $s_0$ e $s_{10}$]
  \label{prop:isolation}
  \begin{enumerate}[label=(\alph*)]
    \item Se $\delta_0 = 0$ e $\delta_{10} = 1$, então
          $\Delta\ks_t = \delta_{20} \xor s_0(x_t)$, logo
          \[s_0(x_t) = \Delta\ks_t \xor \delta_{20}.\]

    \item Se $\delta_0 = 1$ e $\delta_{10} = 0$, então
          $\Delta\ks_t = \delta_{20} \xor s_{10}(x_t)$, logo
          \[s_{10}(x_t) = \Delta\ks_t \xor \delta_{20}.\]
  \end{enumerate}
\end{proposition}

\begin{proof}
  Substituição direta na fórmula do Corolário~\ref{cor:full-diff}. \\
  (a): $\delta_0=0, \delta_{10}=1 \Rightarrow \Delta\ks = \delta_{20} \xor
  (s_0 \AND 1) \xor (0 \AND s_{10}) \xor (0 \AND 1) = \delta_{20} \xor s_0$.\\
  (b): $\delta_0=1, \delta_{10}=0 \Rightarrow \Delta\ks = \delta_{20} \xor
  (s_0 \AND 0) \xor (1 \AND s_{10}) \xor (1 \AND 0) = \delta_{20} \xor s_{10}$. \qedhere
\end{proof}

\subsection{Conjunto de Diferenças Utilizadas}

O código usa $\Delta^{(i)}_0 \in \{x^0, x^1, x^2, x^5, x^{10}\}$ e
propaga cada diferença via $\Delta^{(i)}_t = a^t \Delta^{(i)}_0$.
Para cada clock $t$, verificam-se quais diferenças satisfazem as condições
da Proposição~\ref{prop:isolation}, extraindo $s_0(x_t)$ e $s_{10}(x_t)$.

\begin{theorem}[Equação diferencial geral: verificação]
  O código verifica a equação completa:
  \[
    \Delta\ks_t \stackrel{?}{=}
    \delta_{20} \xor (s_0^* \AND \delta_{10}) \xor (\delta_0 \AND s_{10}^*)
                \xor (\delta_0 \AND \delta_{10}),
  \]
  usando os gabaritos reais $s_0^* = s_0(x_t)$ e $s_{10}^* = s_{10}(x_t)$.
  A taxa de concordância é 1.0 em todas as configurações (prova empírica da
  corretude do Corolário~\ref{cor:full-diff}).
\end{theorem}

% =============================================================================
\section{Módulo 6: \texttt{exhaustive\_speedup}}
\label{sec:speedup}
% =============================================================================

\begin{paperref}
  Biham \& Dunkelman (\S4.1.2, \textit{Improving Exhaustive Key Search of
  Toyocrypt}): \textit{``the key $K \xor s^{126}$ can also be automatically
  discarded without any consideration, reducing the expected time complexity
  of an exhaustive key search by a factor of $2^{126}$ trials. \ldots even
  though the exhaustive key search attack does not require the related-key
  model, it can still benefit from the existence of differentials.''}
  O projeto implementa exatamente este speedup com $\Delta = 1$.
\end{paperref}

\subsection{Propriedade de Simetria dos Pares}

\begin{theorem}[Speedup por complementação via diferencial]
  \label{thm:speedup}
  Seja $\Delta = 1 \in \GF{e}$ com probabilidade-1 característica diferencial
  $\Delta_0 \to \Delta_t = a^t \Delta_0$. Para todo estado inicial $K \in
  \GF{e}$, os fluxos gerados por $K$ e $K' = K \xor \Delta$ satisfazem:
  \[
    \{Z(K), Z(K')\} = \{Z(K'), Z(K)\}
  \]
  no sentido de que os dois fluxos formam um \emph{par identificável} sem
  testar $K'$ de forma independente. A busca exaustiva sobre $2^e$ candidatos
  reduz-se a $2^{e-1}$ testes, um fator 2 de aceleração.
\end{theorem}

\begin{proof}
  Dados os fluxos observados $(Z_{\text{obs}}, Z'_{\text{obs}})$, um
  candidato $K$ satisfaz a hipótese se e somente se um dos dois
  ordenamentos coincide:
  \[
    \bigl[Z(K) = Z_{\text{obs}} \AND Z(K') = Z'_{\text{obs}}\bigr]
    \;\text{OU}\;
    \bigl[Z(K) = Z'_{\text{obs}} \AND Z(K') = Z_{\text{obs}}\bigr].
  \]
  Em ambos os casos, ao encontrar o representante par $K^* = K \AND
  \overline{1}$ (bit 0 zerado), descarta-se automaticamente $K^* \xor 1$.
  Portanto o espaço efetivo de busca é o conjunto de representantes
  $\{K \in \GF{e} : s_0(K) = 0\}$, que tem cardinalidade $2^{e-1}$.
\end{proof}

\subsection{Verificação Experimental}

No código, o speedup teórico de $2\times$ é verificado em $\GF{31}$ com
espaço de $2^{12}$ candidatos: a busca otimizada testa apenas os candidatos
pares ($\Delta=1$ separa par/ímpar), reduzindo exatamente para $2^{11}$
testes. O speedup observado é $2.0\times$ de forma consistente.

% =============================================================================
\section{Módulo 7: \texttt{higher\_order\_stress} e \texttt{higher\_order\_scaling}}
\label{sec:scaling}
% =============================================================================

\begin{paperref}
  Biham \& Dunkelman (\S4, parágrafo sobre monômio de grau 3):
  \textit{``if there is a difference in only one bit which is combined through
  a monomial of degree 3, there is a probability of 1/4 that there is a
  difference in the key streams in the corresponding bit.''}
  Para monômios de grau $d$, a probabilidade cai para $2^{-(d-1)}$, tornando
  o ataque de 1ª ordem inviável. O diferencial de ordem $d$ restaura a
  certeza (probabilidade 1), ao custo de $2^d$ consultas por clock.
\end{paperref}

\subsection{Custo Assintótico do Ataque de Ordem $k$}

\begin{theorem}[Custo mínimo para quebrar grau $d$]
  \label{thm:cost}
  Para uma função filtro $\filter$ de grau algébrico exato $d$, o diferencial
  de ordem mínima que produz uma saída constante (e portanto explorável como
  distinguidor ou para recuperação de estado) é $k = d$, exigindo exatamente
  $2^d$ avaliações de $\filter$ por clock.
\end{theorem}

\begin{proof}
  Segue diretamente do Teorema~\ref{thm:collapse}: $D^d f = \text{const}$ e
  $D^{d-1} f \neq \text{const}$. A definição de $D^k$ envolve $2^k$ pontos
  (subconjuntos de $\{\alpha_1,\ldots,\alpha_k\}$), logo o custo por clock
  é $2^k$. Para $k < d$, o diferencial não é constante e não fornece
  informação utilizável com probabilidade 1. Para $k = d$, o diferencial
  colapsa a uma constante determinística.
\end{proof}

\subsection{Tabela de Viabilidade}

\begin{table}[H]
  \centering
  \caption{Custo do diferencial de ordem $k$ para funções de grau $d = k$}
  \begin{tabular}{crrll}
    \toprule
    Grau $d$ & Ordem $k$ & Consultas/clock & Tempo prático & Viabilidade \\
    \midrule
    2  & 2  & 4            & nanosegundos  & \textbf{Trivial} \\
    3  & 3  & 8            & nanosegundos  & \textbf{Trivial} \\
    5  & 5  & 32           & microssegundos & Fácil \\
    10 & 10 & 1\,024       & milissegundos & Factível \\
    15 & 15 & 32\,768      & segundos      & Custoso \\
    20 & 20 & 1\,048\,576  & horas         & Limite prático \\
    30 & 30 & $\sim 10^9$  & anos          & \textbf{Inviável} \\
    64 & 64 & $\sim 10^{19}$ & $\gg$ idade do universo & \textbf{Seguro} \\
    \bottomrule
  \end{tabular}
\end{table}

\subsection{Análise do \texttt{higher\_order\_stress}}

O módulo compara:
\begin{itemize}
  \item \textbf{[A]} $\filter_{\text{orig}}$ (grau 2): $D^2 \filter_{\text{orig}} = \text{const}$
        $\Rightarrow$ ataque efetivo. Taxa observada: 1.0.
  \item \textbf{[B]} $\filter_{\text{hard}}$ (grau $N-1$, produto AND de $N-1$
        bits): $D^2 \filter_{\text{hard}} \neq \text{const}$
        $\Rightarrow$ ataque de 2ª ordem \textbf{falha}. Taxa observada: $< 1.0$.
\end{itemize}

\begin{proposition}[Função hard: resistência ao $D^2$]
  A função $\filter_{\text{hard}}(x) = \bigwedge_{i=0}^{N-2} s_i(x)$
  tem grau algébrico $N-1$. Portanto $D^2 \filter_{\text{hard}}$ não é
  constante para deltas genéricos, e o ataque de 2ª ordem não a quebra.
\end{proposition}

\begin{proof}
  A ANF de $\filter_{\text{hard}}$ contém o monômio $s_0 s_1 \cdots s_{N-2}$
  de grau $N-1 > 2$. Pelo Teorema~\ref{thm:collapse}, $D^2 \filter_{\text{hard}}$
  ainda contém termos dependentes de $x$ de grau $N-3 > 0$, logo não é constante.
\end{proof}

\subsection{Análise do \texttt{higher\_order\_scaling}}

O módulo implementa a função genérica de grau $G$:
\[
  \filter_G(x) = \bigwedge_{i=0}^{G-1} s_i(x)
                 \xor \bigwedge_{i=0}^{G-2} s_i(x) \xor \cdots \xor s_0(x) \xor s_{G-1}(x),
\]
que é projetada para ter grau algébrico \emph{exato} $G$. O módulo varre
ordens $k = 2, 3, \ldots$ até encontrar o menor $k$ tal que $D^k \filter_G$
é constante (nula) em todos os $t$ testados, confirmando $k = G$.

\begin{theorem}[Verificação empírica do limiar de grau]
  \label{thm:empirical}
  Para $\filter_G$ com grau exato $G$, o módulo \texttt{higher\_order\_scaling}
  detecta $k_{\min} = G$ em todas as configurações testadas
  ($G \in \{2,3,4,5,8,10,15,20\}$), o que constitui evidência empírica
  do Teorema~\ref{thm:collapse}.
\end{theorem}

% =============================================================================
\section{Extensão para Funções Filtro de Grau Maior}
\label{sec:extension}
% =============================================================================

\begin{paperref}
  Biham \& Dunkelman (Sumário, \S7): \textit{``Stream ciphers with high
  probability differentials are susceptible to many attacks: distinguishing
  and key recovery in the related-key/IV model, faster exhaustive key
  searches, and fault analysis.''} Para cifras com funções filtro de alto
  grau, os diferenciais de ordem $d$ ainda se aplicam, porém o custo
  $2^d$ pode tornar o ataque impraticável.
\end{paperref}

\subsection{Generalização do Ataque de Recuperação de Chave}

Para uma função filtro de grau $d$, o Corolário~\ref{cor:d2-filter}
generaliza-se: $D^d f$ produz uma constante que depende apenas das
diferenças $\alpha_1, \ldots, \alpha_d$. Seja:
\[
  C(\alpha_1,\ldots,\alpha_d) = D^d f(x) \quad (\text{independente de } x).
\]
Escolhendo $d$ diferenças tais que $C \neq 0$ (existência garantida pois
o grau é exato $d$), o atacante forma um \emph{distinguidor de ordem $d$}:
observa $2^d$ streams correlacionados $\{Z(x \xor \bigoplus_{i \in S}
\alpha_i)\}_{S \subseteq [d]}$ e verifica se seu XOR é $C$.

\subsection{Generalização da Recuperação de Estado}

Para $d = 2$, o sistema $M \vec{S}_0 = B$ (Seção~\ref{sec:batch}) pode
ser construído porque $\Delta^1 \filter$ é linear em $x_t$. Para $d > 2$:

\begin{itemize}
  \item O diferencial de 1ª ordem $\Delta^1 f$ tem grau $d-1 > 1$, portanto
        não é linearizável diretamente.
  \item O diferencial de ordem $d-1$ produz expressões lineares em $x_t$,
        que podem ser usadas para construir o sistema análogo.
  \item Especificamente, $D^{d-1}_{\alpha_1,\ldots,\alpha_{d-1}} f(x)$ é
        afim em $x$ (grau 1), possibilitando montar uma equação linear por
        clock, com custo $2^{d-1}$ consultas.
  \item O sistema tem $e$ equações (uma por clock) e $e$ incógnitas (bits
        de $x_0$), sendo resolvido por inversão de matriz em $\GFtwo^{e \times e}$
        com $O(e^3)$ operações de bit — idêntico ao caso $d=2$, mas com
        custo de coleta $2^{d-1}$ vezes maior.
\end{itemize}

\subsection{Limiares de Segurança}

\begin{definition}[Grau de segurança de uma função filtro]
  Uma função filtro $\filter$ de grau $d$ é considerada \emph{computacionalmente
  segura} contra ataques diferenciais de ordem $d$ quando:
  \[
    2^d > T_{\text{máx}},
  \]
  onde $T_{\text{máx}}$ é o limiar de consultas tolerado pelo atacante
  por clock de observação.
\end{definition}

Para segurança de 128 bits (padrão moderno), é necessário $d \geq 128$,
o que implica $2^{128} \approx 3.4 \times 10^{38}$ consultas por clock ---
completamente inviável.

\begin{warnbox}
  A função filtro $\filter(x) = s_{20} \xor (s_0 \AND s_{10})$ tem grau
  algébrico $d = 2$, exigindo apenas $2^2 = 4$ consultas por clock para
  o ataque de ordem 2. Isso a torna \textbf{completamente insegura} em
  qualquer instância do projeto, independentemente do tamanho do campo
  $\GF{e}$.
\end{warnbox}

% =============================================================================
\section{Síntese dos Resultados Experimentais}
\label{sec:results}
% =============================================================================

\begin{table}[H]
  \centering
  \caption{Resultados esperados e propriedades matemáticas por módulo}
  \renewcommand{\arraystretch}{1.35}
  \begin{tabular}{lllp{5cm}}
    \toprule
    \textbf{Módulo} & \textbf{Métrica} & \textbf{Resultado} & \textbf{Garantia Matemática} \\
    \midrule
    \texttt{batch\_attack} &
      Acertos / $e$ &
      100\% em todas configs &
      Thm.~\ref{thm:delta-z-s10}: $\Delta\ks_t = s_{10}(x_t)$ c/ prob. 1 \\
    \texttt{recover\_key} &
      Chave recuperada &
      Coincide com semente &
      Thm.~\ref{thm:invertible}: $M$ invertível sobre $\GFtwo$ \\
    \texttt{diff\_corr} &
      $\phi(\Delta\ks, s_{10})$ &
      1.0 &
      Thm.~\ref{thm:diff-corr}: correlação perfeita \\
    \texttt{dist\_attack} &
      Taxa de discordância &
      0\% (alvo), ~50\% (rand.) &
      Thm.~\ref{thm:distinguisher}: $\delta_0=\delta_{10}=0 \Rightarrow \Delta\ks=\delta_{20}$ \\
    \texttt{higher\_ord.} &
      Taxa $D^2 = \text{pred}$ &
      1.0 &
      Cor.~\ref{cor:d2-filter}: $D^2 f = \text{const}$ para grau 2 \\
    \texttt{multi\_diff} &
      Equações corretas &
      100\% &
      Cor.~\ref{cor:full-diff}: equação diferencial completa \\
    \texttt{exh\_speedup} &
      Fator de aceleração &
      $2.0\times$ &
      Thm.~\ref{thm:speedup}: pares $\{K, K\xor\Delta\}$ eliminados \\
    \texttt{stress\_test} &
      $D^2$ p/ grau-$(N-1)$ &
      Inconsistente &
      Thm.~\ref{thm:collapse}: $D^{d-1} f \neq \text{const}$ \\
    \texttt{scaling} &
      $k_{\min}$ para grau $G$ &
      $k_{\min} = G$ &
      Thm.~\ref{thm:empirical}: limiar exato do colapso \\
    \bottomrule
  \end{tabular}
\end{table}

% =============================================================================
\section{Conclusões}
\label{sec:conclusion}
% =============================================================================

Este relatório demonstrou matematicamente que o projeto implementa de forma
rigorosa e completa o framework de criptanálise diferencial de Biham \&
Dunkelman aplicado a cifras de fluxo baseadas em LCG sobre $\GF{e}$.

Os resultados centrais provados são:

\begin{enumerate}[label=(\roman*)]
  \item \textbf{Propagação determinística (Thm.~\ref{thm:linear-prop}):}
        $\Delta_t = a^t \Delta_0$ — a linearidade do LCG garante que o
        atacante conhece a diferença interna em qualquer instante.

  \item \textbf{Extração de bit por diferença (Thm.~\ref{thm:delta-z-s10}):}
        $\filter(x) \xor \filter(x \xor 1) = s_{10}(x)$ — a função filtro
        de grau 2 lineariza-se sob diferença no bit 0.

  \item \textbf{Recuperação de chave por sistema linear (Thm.~\ref{thm:invertible}):}
        a matriz $M \in \GFtwo^{e \times e}$ é invertível, e $x_0 = M^{-1}B$.

  \item \textbf{Colapso do diferencial de ordem $d$ (Thm.~\ref{thm:collapse}):}
        $D^d f = \text{const}$ é a condição necessária e suficiente para a
        exploração, exigindo $2^d$ consultas por clock.

  \item \textbf{Inviabilidade para grau alto (Thm.~\ref{thm:cost}):}
        para $d \geq 30$, o custo $2^d > 10^9$ consultas/clock torna o ataque
        impraticável; para $d \geq 64$, é computacionalmente seguro.
\end{enumerate}

A vulnerabilidade fundamental da cifra reside no \textbf{baixo grau algébrico
($d=2$) da função filtro}, que permite todos os ataques descritos com custo
trivial (4 consultas/clock). Aumentar $d$ para $\geq 30$ eliminaria os ataques
diferenciais práticos, mas introduziria outras vulnerabilidades (não-linearidade
insuficiente, correlação com sub-funções, etc.) que requerem análise adicional.

% =============================================================================
\appendix
% =============================================================================

\section{Identidades de Álgebra Booleana Utilizadas}
\label{app:boolean}

As seguintes identidades sobre $(\GFtwo, \xor, \AND)$ são usadas nas provas:

\begin{align*}
  a \xor a &= 0 \tag{A1: Complemento} \\
  a \xor 0 &= a \tag{A2: Identidade XOR} \\
  a \AND (b \xor c) &= (a \AND b) \xor (a \AND c) \tag{A3: Distributividade} \\
  a \AND 0 &= 0, \quad a \AND 1 = a \tag{A4: Absorção} \\
  a \xor b &= b \xor a \tag{A5: Comutatividade} \\
  (a \xor b) \xor c &= a \xor (b \xor c) \tag{A6: Associatividade}
\end{align*}

\section{Propriedades de $\GF{e}$ Utilizadas}
\label{app:gf2e}

\begin{enumerate}[label=(GF\arabic*)]
  \item $\GF{e} = \GFtwo[x]/\langle h(x) \rangle$ com $h$ primitivo de grau $e$.
  \item A multiplicação por $a \in \GF{e}$ é um automorfismo $\GFtwo$-linear.
  \item O grupo multiplicativo $\GF{e}^*$ é cíclico de ordem $2^e - 1$.
  \item A projeção $s_k : \GF{e} \to \GFtwo$, $s_k(y) = [x^k] y$, é
        $\GFtwo$-linear e não-nula.
  \item A base $\{\xi^0, \xi^1, \ldots, \xi^{e-1}\}$ é uma base de $\GF{e}$
        como $\GFtwo$-espaço vetorial.
\end{enumerate}

\section{Referência: Paper Original}
\label{app:paper}

\begin{itemize}
  \item \textbf{Autores:} Eli Biham (Technion, Haifa, Israel) e
        Orr Dunkelman (KU Leuven, Bélgica).
  \item \textbf{Título:} \textit{Differential Cryptanalysis in Stream Ciphers}.
  \item \textbf{Arquivo:} \texttt{docs/Differential Cryptanalysis in Stream Ciphers.pdf}.
  \item \textbf{Cifras analisadas no paper:} Toyocrypt-HS1, A5/1, RC4,
        Trivium, LILI-128, LILI-II, Py/PyPy, Helix.
  \item \textbf{Correspondência com o projeto:}
    \begin{itemize}
      \item LCG $x_{t+1} = ax_t + c$ em $\GF{e}$ $\leftrightarrow$ LFSR
            primitivo de $e$ estágios.
      \item $\filter(x) = s_{20} \xor s_0 s_{10}$ $\leftrightarrow$
            função combinadora de Toyocrypt-HS1 (simplificada a 3 variáveis).
      \item $\Delta_0 = 1$ $\leftrightarrow$ diferença $\Delta S = s_0$
            de Biham \& Dunkelman (\S4.1.3).
    \end{itemize}
\end{itemize}

\end{document}
